% Revize: vyhradne narocnejsi soustavy 3x3, navaznost na 83.
% Priklady z ucitelem dodane sbirky 31b,31c,31d,31e; metodika 020313-020314.
\documentclass[11pt,a4paper]{article}
\input{../../sablony/pripravy/priprava.sty}
\nastavPripravu{1. ročník}{Rovnice a nerovnice}{84}{Soustavy tří rovnic o třech neznámých}
\begin{document}
\blok{Návaznost a záměr}
Z minulé hodiny známe první eliminační krok. Dnes řešíme samostatnější, algebraicky náročnější soustavy $3\times3$. Cílem není hádat tři čísla, ale \emph{systematicky} převést soustavu $3\times3$ na $2\times2$, poté na $1\times1$ a dosazovat zpět. Stále provádíme pouze úpravy lineárních rovnic.
\blok{Přehled metody}
\begin{enumerate}[label=\arabic*.]
\item Zvol jednu neznámou k eliminaci (ideálně s koeficientem $1$ nebo $-1$).
\item Vytvoř \textbf{dvě nové rovnice} o zbývajících dvou neznámých.
\item Vyřeš jejich soustavu jednou známou metodou a dopočti třetí neznámou.
\item Proveď zkoušku ve všech \textbf{třech původních} rovnicích.
\end{enumerate}
\begin{prikladReseniBox}[title={Příklad 1 – sbírka č. 31b}]
$\begin{cases}2x+3y+z=15,\\7x-y+z=9,\\x+2y+z=9.\end{cases}$
\tcblower
Odečteme třetí rovnici od první a od druhé:
\begin{align*}
(\mathrm I)-(\mathrm{III})&:\quad x+y=6,\\
(\mathrm{II})-(\mathrm{III})&:\quad 6x-3y=0,\\
&\phantom{:}\quad y=2x,\\
x+2x&=6,\\
x&=2,\qquad y=4,\\
x+2y+z&=9,\\
2+8+z&=9,\\
z&=-1.
\end{align*}
Zkouška: $4+12-1=15$, $14-4-1=9$, $2+8-1=9$. $K=\{[2;4;-1]\}$.
\end{prikladReseniBox}
\begin{prikladReseniBox}[title={Příklad 2 – sbírka č. 31d, volba neznámé}]
$\begin{cases}2x-y=6,\\y+4z=8,\\x-z=1.\end{cases}$
\tcblower
\begin{align*}
x&=z+1,\\
y&=2x-6=2z-4,\\
(2z-4)+4z&=8,\\
6z&=12,\\
z&=2,\\
x&=3,\qquad y=0.
\end{align*}
Zkouška všech tří rovnic: $6-0=6$, $0+8=8$, $3-2=1$.
\end{prikladReseniBox}
\poznamka{Pokud je třída pomalejší, druhý společný příklad využij jako rezervu. Povinné jádro je zvládnutí jednoho úplného eliminačního postupu a jeho zkouška.}
\blok{Samostatné procvičování}
\begin{zadaniReseniBox}[title={Úloha ze sbírky č. 31c – hlídej si pořadí souřadnic}]
$\begin{cases}2x+y-z=0,\\4x+2y+z=0,\\x-y+3z=0.\end{cases}$
\tcblower
\begin{align*}
(\mathrm{II})-2(\mathrm I)&:\quad3z=0,\\
z&=0,\\
2x+y&=0,\\
x-y&=0,\\
y&=x,\\
3x&=0\ \Rightarrow\ x=0,\quad y=0.
\end{align*}
$K=\{[0;0;0]\}$. Nulová pravá strana \emph{neznamená automaticky} nekonečně mnoho řešení.
\end{zadaniReseniBox}
\blok{Volitelná rezerva – může nastat spor?}
\begin{zadaniReseniBox}[title={Úloha ze sbírky č. 31e}]
$\begin{cases}2x+y-z=0,\\x+y+2z=4,\\4x+3y+3z=5.\end{cases}$
\tcblower
\begin{align*}
y&=-2x+z,\\
-x+3z&=4\ \Rightarrow\ x=3z-4,\\
-2x+6z&=5,\\
-2(3z-4)+6z&=5,\\
8&=5,\\
K&=\varnothing.
\end{align*}
\end{zadaniReseniBox}
\blok{Domácí úkol}
% MAT-REV-2026-11-07: homework-84
\begin{zadaniReseniBox}[title={Domácí příprava}]
\begin{enumerate}[label=\alph*)]
\item $\begin{cases}3x-2y+z=11,\\2x+y-3z=-6,\\4x+3y+2z=11.\end{cases}$
\item $\begin{cases}3x+2y+z=3,\\x+y+z=2,\\4x+3y+2z=5.\end{cases}$
\item $\begin{cases}x+y+z=2,\\2x-y+3z=8,\\3x+2y-z=3,\end{cases}$.
\item $\begin{cases}x+y+z=1,\\2x-y+z=-5,\\x-2y+3z=-12,\end{cases}$.
\item $\begin{cases}x+y+z=4,\\2x+2y+2z=8,\\x-y+z=0,\end{cases}$.
\end{enumerate}
\tcblower
\begin{enumerate}[label=\alph*)]
\item
Z druhé rovnice $y=-6-2x+3z$. Dosadíme do první a třetí:
\[\begin{aligned}
7x-5z&=-1,\\
-2x+11z&=29,\\
2(7x-5z)+7(-2x+11z)&=-2+203,\\
67z&=201,\\
z&=3,\\
7x-15&=-1\ \Rightarrow\ x=2,\\
y&=-6-4+9=-1.
\end{aligned}\]
$K=\{[2;-1;3]\}$. Zkouška ve všech původních rovnicích.

\item
Třetí rovnice je součtem prvních dvou, neklade novou podmínku:
\[\begin{aligned}
(\mathrm I)-2(\mathrm{II})&:\quad x-z=-1,\\
x&=z-1,\\
y&=2-x-z=3-2z.
\end{aligned}\]
Položíme $z=t$ a dostaneme $K=\{[t-1;3-2t;t]\mid t\in\mathbb R\}$.

\item
\[\begin{aligned}
x&=2-y-z\\-3y+z&=4\\-y-4z&=-3\\y&=3-4z\\13z&=13\\z&=1,\quad y=-1,\quad x=2\\K&=\{[2;-1;1]\}
\end{aligned}\]

\item
\[\begin{aligned}
x&=1-y-z\\3y+z&=7\\3y-2z&=13\\3z&=-6\\z&=-2,\quad y=3,\quad x=0\\K&=\{[0;3;-2]\}
\end{aligned}\]

\item
\[\begin{aligned}
x+y+z&=4,\quad x-y+z=0\\2y&=4\ \Rightarrow\ y=2\\x+z&=2\\z&=t,\quad x=2-t\\K&=\{[2-t;2;t]\mid t\in\mathbb R\}
\end{aligned}\]
\end{enumerate}
\end{zadaniReseniBox}


\poznamka{Čas: návrat k eliminaci 4 min, příklad 1 přibližně 15 min, příklad 2 asi 10 min, samostatná práce 12 min, shrnutí 4 min. Rezervu nech až po samostatném úspěchu.}
\end{document}
